\section{Related Work}
\label{sec:related}

Table~\ref{tab:related} compares compiler and incremental systems by their
programmable units, composition boundaries, and dependency mechanisms.

\begin{table*}[t]
  \caption{Programmable units and dependency mechanisms. The two bands share
  comparison dimensions; entries name mechanisms rather than capability scores.}
  \label{tab:related}
  \normalsize
  \setlength{\tabcolsep}{4pt}
  \renewcommand{\arraystretch}{1.08}
  \begin{tabularx}{\textwidth}{@{}>{\raggedright\arraybackslash}p{0.14\textwidth}*{5}{>{\raggedright\arraybackslash}X}@{}}
    \toprule
    & MLIR\newline\cite{lattner2021mlir,mlirpass}
    & xDSL\newline\cite{fehr2025xdsl}
    & TVM\newline\cite{chen2018tvm,feng2022tensorir}
    & Exo 2\newline\cite{ikarashi2025exo2} & \textbf{Joggle} \\
    \midrule
    Programming & C++ / ODS & Python & Python / C++ & Python & \textbf{Typed jog calls} \\
    Composition & Dialects, passes & Dialects, passes & Tensor programs, schedules & Scheduling libraries & \textbf{Cross-stage mods} \\
    References & Operations, values & Operations, values & Tensor blocks & Cursors & \textbf{Typed graph handles} \\
    Dependency mechanism & Analysis preservation & SSA use-def links & Dataflow, schedules & Cursor forwarding & \textbf{Reads, effects, revisions} \\
    \midrule
    & Transform\newline\cite{lucke2025transform}
    & \texttt{egg}\newline\cite{willsey2021egg}
    & rustc\newline\cite{rustcincremental}
    & Adapton\newline\cite{hammer2014adapton}
    & Build systems\newline\cite{mokhov2018build} \\
    \midrule
    Programming & Transform IR & Rust rules & Rust queries & Host-language API & Task rules \\
    Composition & Transform operations & Rewrite rules & Query calls & Thunk calls & Build rules \\
    References & Payload handles & E-classes & Query keys & Thunks & Task keys \\
    Dependency mechanism & Handle effects & E-class rebuilding & Red-green query DAG & Demanded computation graph & Task dependency graph \\
    \bottomrule
  \end{tabularx}
\end{table*}

\subsection{Extensible Compiler Infrastructures}

LLVM supplies typed SSA, analyses, and passes~\cite{lattner2004llvm}; MLIR adds
extensible dialects and multi-level lowering~\cite{lattner2021mlir}. IRDL specifies
operations, types, attributes, and constraints declaratively~\cite{fehr2022irdl},
while xDSL brings MLIR-compatible construction into Python~\cite{fehr2025xdsl}.
MLIR's pass manager caches analyses at operation anchors and invalidates them
according to preservation declarations~\cite{mlirpass}.

Metaprogramming offers another route to extensibility. LMS stages code through
types~\cite{rompf2010lms}; AnyDSL specializes higher-order
functions through partial evaluation~\cite{leissa2018anydsl}. These approaches
make program generation compositional. Our compiler functions instead expose
the roles acting on a mutable graph through one call and value model. Mods
give those functions a shared ownership and publication boundary, orthogonal
to program containment.

\subsection{Tensor Compilers}

Halide separates algorithms from schedules~\cite{ragankelley2012halide}. TVM
combines graph and tensor optimization~\cite{chen2018tvm}, TensorIR exposes tensor
computation blocks~\cite{feng2022tensorir}, and Ansor searches a hierarchical
space of tensor programs~\cite{zheng2020ansor}. Triton exposes tiled
kernels~\cite{tillet2019triton}. At graph level, TASO generates and verifies
substitutions~\cite{jia2019taso}, while Mirage searches transformations across
kernel, thread-block, and thread levels~\cite{wu2025mirage}.

PluS packages expert subgraph optimizations as pluggable graph
schedules~\cite{wu2025plus}. ONNX-MLIR lowers model operations~\cite{jin2020onnxmlir},
and ONNX Runtime assigns graph partitions to execution providers~\cite{ortarchitecture}.
These are complementary optimization and deployment boundaries. A mod instead
groups the definitions, analyses, rewrites, and emitters implementing one
cross-stage feature. Section~\ref{sec:evaluation} measures extension changes
separately from the performance of generated artifacts.

\subsection{Programmable Transformations}

Nanopass derives representation checks and traversal support from explicit
source and target languages~\cite{keep2013nanopass}. Exo externalizes accelerator
instructions and scheduling decisions~\cite{ikarashi2022exo}; Exo 2 adds reusable
scheduling libraries through actions, inspection, and cursors~\cite{ikarashi2025exo2}.
The Transform dialect represents schedules as IR with payload handles and
effects~\cite{lucke2025transform}. For rule-driven optimization, \texttt{egg}
combines equality saturation and e-class analyses~\cite{willsey2021egg};
\texttt{egglog} integrates these with Datalog and incremental
evaluation~\cite{zhang2023egglog}.

For agents, this broader extension surface moves beyond fast-kernel
generation~\cite{ouyang2025kernelbench} to matched tasks across six compiler
roles, all with native build-and-test feedback.

\subsection{Incremental Computation}

Self-adjusting computation combines dynamic dependence graphs and memoization
to reuse executions after input changes~\cite{acar2009selfadjusting}. Adapton
adds demand-driven composition~\cite{hammer2014adapton}. For program analysis,
IncA compiles specifications into incrementally maintained graph
patterns~\cite{szabo2016inca}. Build Systems à la Carte separates dependency
discovery, scheduling, and rebuilding~\cite{mokhov2018build}; rustc reuses
validated query results through a red-green dependency graph~\cite{rustcincremental}.
LLVM ORC instead manages on-demand materialization and symbol dependencies
for JIT compilation~\cite{llvmorc}.

For mutable compiler graphs, our evaluator records typed observations and
changed scopes. Revisions validate state, generations validate identity, and
effect overlap selects downstream work. Transactions publish graph and
dependency records together.
